\chapter{Introduction}
\label{ch:intro}

\section{The Classical Minimal Surface Problem}

\subsection{Historical Origins: From Plateau to Morse Theory}

The study of minimal surfaces has its roots in eighteenth-century physics and the calculus of
variations. The name ``minimal surface'' itself originates from the physical observation that soap
films spanning a wire frame naturally minimize area subject to the boundary constraint. This
physical phenomenon, formalized as the \emph{Plateau problem}, asks: given a Jordan curve
$\gamma \subset \mathbb{R}^3$, does there exist a surface of least area with boundary $\gamma$?

The Plateau problem was solved in the 1930s independently by Douglas and Rad\'o using conformal
parametrizations and harmonic maps. Their work established that smooth minimal surfaces exist for
nice boundary data, opening the door to a systematic study of minimal submanifolds in Riemannian
geometry.

A fundamentally different approach emerged in the 1960s through the work of Birkhoff on closed
geodesics. Birkhoff's \emph{min-max method} constructs a nontrivial closed geodesic on the
two-sphere $S^2$ by considering a one-parameter family of curves (a ``sweepout'') connecting a
point curve to itself with nonzero degree. The minimax principle guarantees that the curve of
maximal length in this family can be shortened to a closed geodesic. This variational idea---finding
critical points by analyzing the maximum value of a functional over continuous families---became
the foundation of modern min-max theory.

The extension of Birkhoff's method to higher-dimensional minimal submanifolds, however, faced
formidable obstacles. Unlike curves in surfaces (where regularity is almost automatic), minimal
hypersurfaces in higher dimensions may develop singularities, and controlling the topology
(embeddedness, genus bounds) requires sophisticated machinery from geometric measure theory. These
challenges motivated the development of \emph{Almgren--Pitts theory}, the modern framework for
min-max constructions of minimal submanifolds.

\subsection{The Three Fundamental Questions}

The classical theory of minimal submanifolds revolves around three fundamental questions that have
shaped the development of the field:

\begin{question}[Existence]
Given topological or boundary constraints, does there exist a minimal submanifold satisfying these
conditions?
\end{question}

The existence question has been addressed through various approaches: direct methods in the
calculus of variations (minimization), geometric flows (mean curvature flow, Ricci flow), and
min-max methods. Each approach has advantages and limitations. Direct minimization often produces
area-minimizers but may fail to capture non-minimizing critical points. Geometric flows can lose
topology through singularities. Min-max methods, while technically demanding, have proven remarkably
effective at constructing minimal submanifolds with prescribed topological features.

\begin{question}[Regularity]
Are the minimal submanifolds produced by variational methods smooth? If singularities occur, what
is their structure?
\end{question}

The regularity question lies at the heart of geometric analysis. Even when existence is established
through variational or topological methods, the resulting objects may be singular (as in geometric
measure theory) or only weakly regular. Understanding when and why minimal submanifolds are smooth
requires deep techniques from partial differential equations, harmonic analysis, and the theory of
elliptic regularity. As we will see in Section~\ref{sec:regularity}, the answer depends critically
on dimension: there is a fundamental \emph{dimension barrier} beyond which singularities are
unavoidable.

\begin{question}[Multiplicity and Structure]
How many minimal submanifolds exist? How are they distributed in the ambient manifold?
\end{question}

The multiplicity question, exemplified by Yau's conjecture (existence of infinitely many minimal
hypersurfaces in closed manifolds), asks not just whether minimal submanifolds exist, but whether
they are abundant. Recent breakthroughs have shown that in generic closed manifolds, minimal
hypersurfaces are not only infinite in number, but form a \emph{dense} set and exhibit remarkable
\emph{equidistribution} properties with respect to the ambient volume measure. These results reveal
deep connections between minimal surface theory, spectral geometry, and dynamical systems.

\section{Min-Max Theory: From Existence to Multiplicity}\label{sec:minmax}

Having established the three fundamental questions, we now turn to \emph{min-max theory}, the
primary tool for addressing existence and multiplicity in the closed manifold setting. Min-max
methods have achieved spectacular successes over the past decade, resolving long-standing
conjectures and revealing unexpected structure in the space of minimal submanifolds.

\subsection{The Almgren--Pitts Framework}

The modern min-max theory for minimal hypersurfaces descends from the pioneering work of Almgren
\cite{Alm62,Alm65,Alm68} and Pitts \cite{Pit81}. Pitts adapted Almgren's homotopy-theoretic
framework to the setting of codimension-one varifolds in Riemannian manifolds, producing the first
general existence theorem for minimal hypersurfaces via min-max methods.

\subsubsection{Sweepouts and Width Functionals}

The basic setup is as follows. Let $(M^{n+1}, g)$ be a closed Riemannian manifold. A
\emph{sweepout} is a continuous map
\[
    \Phi: X \to \mathcal{Z}_n(M; \mathbb{Z}_2)
\]
from a parameter space $X$ (typically a sphere or torus) to the space $\mathcal{Z}_n(M;\mathbb{Z}_2)$
of $n$-dimensional mod-$2$ flat chains, satisfying certain non-triviality conditions encoded by
Almgren's cohomology class $\lambda_n \in H^1(\mathcal{Z}_n(M;\mathbb{Z}_2); \mathbb{Z}_2)$.

For each sweepout $\Phi$, the \emph{width} is defined as
\[
    \mathbf{L}(\Phi) := \max_{x \in X} \mathcal{H}^n(\Phi(x)),
\]
where $\mathcal{H}^n$ denotes $n$-dimensional Hausdorff measure. The \emph{min-max value} (or
\emph{$1$-width}) is then
\[
    \omega_1(M,g) := \inf_{\Phi \in \Pi} \mathbf{L}(\Phi),
\]
where $\Pi$ denotes a suitable homotopy class of sweepouts.

Pitts' fundamental theorem asserts that this min-max value is realized by a stationary integral
varifold---a weak notion of minimal hypersurface that may have singularities and integer
multiplicities. The proof involves a delicate pull-tight argument and uses compactness of varifolds
to extract a limiting object from a minimizing sequence of sweepouts.

\subsubsection{Continuous Sweepouts: The Colding--De Lellis Simplification}

A decisive simplification of the Almgren--Pitts theory was achieved by Colding--De Lellis \cite{CL03}
and De Lellis--Tasnady \cite{DLT13}. They introduced \emph{continuous sweepouts}---families of
hypersurfaces parametrized continuously without discretization---and developed a streamlined
pull-tight procedure that clarifies the geometric picture.

The continuous sweepout framework has several advantages: it is more intuitive, easier to adapt to
settings with boundary, and better suited for tracking geometric quantities (such as relative
boundaries in the free-boundary setting). This approach forms the technical foundation for the main
result of this thesis (Theorem~\ref{thm:main}).

\subsection{The Weyl Law and Volume Spectrum (Codimension One)}

A major breakthrough in min-max theory was the proof of the \emph{Weyl law for the volume spectrum}
by Liokumovich--Marques--Neves \cite{LMN}. This result provides quantitative control over the growth
rate of higher min-max values and has far-reaching consequences for multiplicity questions.

\subsubsection{The Volume Spectrum}

For $p \geq 1$, the \emph{$p$-width} $\omega_p^{n-1}(M,g)$ is defined by considering
$p$-parameter sweepouts---families of hypersurfaces parametrized by a $p$-dimensional parameter
space $X$ satisfying a non-triviality condition encoded by the $p$-th cup power $\lambda_{n-1}^p$
of Almgren's cohomology class. The sequence $\{\omega_p^{n-1}(M,g)\}_{p=1}^\infty$ is called the
\emph{volume spectrum}.

In the 1980s, Gromov \cite{Gro86,Gro02,Gro09} proposed viewing the volume spectrum as a nonlinear
analog of the eigenvalues of the Laplacian, and conjectured an asymptotic \emph{Weyl law}:
\[
    \lim_{p \to \infty} \omega_p^{n-1}(M,g) \, p^{-\frac{n-1}{n}} = \alpha(n) \, \mathrm{Vol}(M,g)^{\frac{n-1}{n}}
\]
for a universal constant $\alpha(n) > 0$ depending only on the dimension $n$.

\subsubsection{Resolution of the Weyl Law}

Liokumovich--Marques--Neves \cite{LMN} proved Gromov's conjecture for hypersurfaces ($k = n-1$),
establishing the Weyl law with sharp asymptotics. Their proof combines the Almgren--Pitts min-max
construction with sophisticated regularity arguments and a delicate analysis of how widths vary
under conformal deformations.

For higher codimension, Gromov's conjecture for $1$-cycles (geodesic nets) was first proved by
Guth--Liokumovich \cite{GL22} in dimension $n=3$ using parametric versions of the isoperimetric
and coarea inequalities, and then by Staffa \cite{Sta-Weyl,Sta-Coarea} in arbitrary dimension $n$,
resolving the conjecture for codimension $n-1$.

The Weyl law has profound implications:
\begin{itemize}
    \item It implies the existence of \emph{infinitely many} distinct min-max minimal hypersurfaces
    (since $\omega_p$ grows without bound as $p \to \infty$).

    \item It provides quantitative estimates on the areas of these hypersurfaces.

    \item It establishes a deep analogy between the volume spectrum and the Laplace spectrum,
    suggesting that minimal hypersurfaces encode spectral information about the ambient manifold.
\end{itemize}

\subsection{Yau's Conjecture and Recent Breakthroughs}

Building on the Weyl law, a series of spectacular results resolved Yau's conjecture and revealed
unexpected structure in the space of minimal hypersurfaces.

\subsubsection{Infinitely Many Minimal Hypersurfaces}

In the early 1980s, Yau conjectured that every closed Riemannian three-manifold contains infinitely
many immersed minimal surfaces. This conjecture was first resolved by Marques--Neves \cite{MN17}
for manifolds with positive Ricci curvature in dimensions up to seven, using the Weyl law combined
with a careful analysis of multiplicities.

Subsequently, Song \cite{Song18} proved Yau's conjecture for arbitrary closed Riemannian manifolds
of dimension three to seven, removing the Ricci curvature assumption. Song's breakthrough relied on
a delicate Lusternik--Schnirelmann type argument to handle the \emph{multiplicity problem}: the
same minimal hypersurface appearing with different integer multiplicities represents distinct
critical points from the variational perspective, and these must be carefully distinguished to
ensure infinitude of geometrically distinct hypersurfaces.

\subsubsection{Generic Density}

The progression from existence to geometric understanding continued with Irie--Marques--Neves
\cite{IMN17}, who proved a \emph{generic density theorem}: for Baire-generic metrics on closed
manifolds (dimension three to seven), the union of all closed embedded minimal hypersurfaces is
\emph{dense} in $M$.

This result settled Yau's conjecture in the generic case and revealed a much stronger structure
than mere infinitude: the minimal hypersurfaces collectively fill the ambient space in a
topological sense. The proof uses bumpy metrics (metrics for which all minimal hypersurfaces are
non-degenerate) combined with Morse-theoretic arguments and the Weyl law.

\subsubsection{Generic Equidistribution}

Shortly thereafter, Marques--Neves--Song \cite{MNS} refined the density result into a quantitative
\emph{generic equidistribution theorem}: for generic metrics, there exists a sequence
$\{\Sigma_j\}$ of closed embedded minimal hypersurfaces such that
\[
    \lim_{q \to \infty} \frac{1}{\sum_{j=1}^q \vol(\Sigma_j)} \sum_{j=1}^q \int_{\Sigma_j} f \, d\mu
    = \frac{1}{\vol(M)} \int_M f \, d\vol
\]
for any smooth function $f: M \to \mathbb{R}$.

This equidistribution result demonstrates that minimal hypersurfaces collectively encode the
ambient geometry in a profound way, extending classical equidistribution phenomena from dynamical
systems and number theory to geometric variational problems. The proof crucially relies on the
Weyl law and quantitative control over how the $p$-widths vary under conformal deformations of the
metric. A parallel equidistribution result for stationary geodesic nets in $n$-manifolds ($n \geq 3$)
was established by Li--Staffa \cite{LS23}, building on Staffa's Weyl law and Liokumovich--Staffa's
generic density result \cite{LioStaffa-dense}.

\section{Regularity Theory: The Dimension Barrier and Recent Breakthroughs}\label{sec:regularity}

Having discussed existence and multiplicity (Question 1 and Question 3 from
Section~\ref{sec:minmax}), we now turn to the \emph{regularity question}: Are minimal hypersurfaces
produced by variational methods smooth? The answer to this question depends critically on the
dimension of the ambient manifold and reveals deep connections between geometric analysis, partial
differential equations, and the calculus of variations.

\subsection{Classical Regularity Theory and the Dimension Barrier}

The regularity theory for minimal hypersurfaces has a rich history dating back to the early 1970s.
A fundamental result of Federer \cite{Fed70}, building on Simons' pioneering work \cite{Sim68} on
the nonexistence of stable cones in $\mathbb{R}^{n+1}$ for $2 \leq n \leq 6$, established that
area-minimizing rectifiable hypercurrents have singular sets of Hausdorff codimension at least $7$
away from their boundary. In particular, for $2 \leq n \leq 6$, any area-minimizing integral
$n$-cycle in a Riemannian manifold $M^{n+1}$ must be a smooth minimal hypersurface.

However, the situation changes dramatically in dimension $n = 7$. Bombieri--De Giorgi--Giusti
\cite{BDGG69} proved that the \emph{Simons cone}
\[
    \mathcal{C} = \{(x,y) \in \mathbb{R}^4 \times \mathbb{R}^4 : |x| = |y|\}
\]
is area-minimizing in $\mathbb{R}^8$, yet has an isolated singularity at the origin. This landmark
example establishes dimension $8$ as the first dimension where singularities can appear for
area-minimizing hypersurfaces, creating a fundamental \emph{dimension barrier} for regularity theory.

The existence of the Simons cone has profound implications: it shows that even the strongest
variational condition (area-minimization) does not guarantee smoothness in dimensions $\geq 8$,
and that singularities are an intrinsic feature of the geometry in higher dimensions.

\subsection{Integrality and Wickramasekera's Regularity Theorem}\label{subsec:regularity-stable}

For minimal hypersurfaces in \emph{closed manifolds} arising from variational methods (such as
min-max theory), the relevant notion is \emph{stability} rather than area-minimization. This
regularity theory, culminating in Wickramasekera's theorem, forms the theoretical foundation for
Result II of this thesis (Theorem~\ref{thm:non-excessive}). A minimal hypersurface $\Sigma$ is
\emph{stable} if the second variation of area is non-negative for all compactly supported normal
variations. In the language of varifolds, this translates to the stability inequality
\[
    \int_{\text{spt}\|V\|} |\nabla \phi|^2 \, d\|V\|
    \geq \int_{\text{spt}\|V\|} |A|^2 \phi^2 \, d\|V\|
\]
for all compactly supported smooth functions $\phi$, where $A$ denotes the second fundamental form.

The regularity theory in this setting was pioneered by Schoen--Simon \cite{SS81} and
Schoen--Simon--Yau \cite{SSY75}, who proved that stable minimal hypersurfaces with natural a priori
assumptions on the singular set are smooth in dimensions up to $7$. Their work extended Simons'
nonexistence results for stable cones and applied deep PDE techniques to control the regularity
near potential singularities.

A decisive breakthrough came with Wickramasekera's deep regularity theorem \cite{Wic14}, which
provides optimal regularity for stable minimal hypersurfaces (varifolds) under the crucial
assumption of \emph{integrality}.

\begin{definition}[Integral Varifold]
An integral varifold is one whose tangent planes have integer multiplicities almost everywhere.
\end{definition}

Wickramasekera's theorem asserts that if $V$ is an integral stationary varifold in $M^{n+1}$ with
$2 \leq n \leq 6$ satisfying the stability inequality, then $V$ is a smooth minimal hypersurface
(possibly with integer multiplicity).

\paragraph{Integrality: The Crucial Bridge for Closed Manifolds.}
For the closed manifold setting, integrality is the \emph{crucial bridge} between variational
min-max constructions and regularity theory. Without integrality, Wickramasekera's regularity
theorem cannot be applied, and smoothness cannot be guaranteed. In the classical Almgren--Pitts
framework, integrality is built into the construction through the use of integral currents and
careful pull-tight arguments. However, in more modern frameworks (such as the non-excessive approach
discussed in Section~\ref{sec:free-bdry}), integrality is not automatic and must be established
separately, creating a fundamental challenge that we address in Result II of this thesis (see
Theorem~\ref{thm:non-excessive}).

The quantitative refinements and rectifiability of singular sets have been further studied by
Cheeger--Naber \cite{CN13} and Naber--Valtorta \cite{NV20}, providing sharp estimates on the
structure of singularities when they do occur.

\subsection{Regularity for Free Boundary Minimal Hypersurfaces}\label{subsec:lizhou-regularity}

The free boundary setting requires specialized regularity theory that differs fundamentally from
the closed case. Li--Zhou \cite{LZ17} developed a comprehensive regularity framework for free
boundary minimal hypersurfaces (FBMH) arising from min-max constructions, establishing smoothness
up to the boundary and verifying the orthogonality condition $\nu_\Sigma \perp T(\partial M)$ in
dimensions $3 \leq n+1 \leq 7$.

Their approach exploits the \emph{doubling trick}: by reflecting the manifold $M$ across its
boundary $\partial M$, one obtains a doubled manifold $\tilde{M}$ without boundary, allowing the
use of interior regularity techniques. The key technical challenge is to ensure that the reflected
varifold inherits appropriate regularity properties. Li--Zhou establish this through an
\emph{almost minimizing property in annuli intersecting the boundary}, which provides sufficient
control to apply blow-up analysis and conclude smoothness.

Crucially, the Li--Zhou regularity framework operates directly within the min-max construction and
does not require a separate integrality verification step as in the Wickramasekera approach for
closed manifolds. The free boundary orthogonality condition emerges naturally from the variational
principle and the reflection symmetry.

This regularity theory forms the foundation for Result I of this thesis (Theorem~\ref{thm:main}),
allowing us to construct smooth free boundary minimal hypersurfaces in complete manifolds with
boundary.

\subsection{Generic Regularity: From Dimension 8 to Dimension 11}\label{subsec:generic-regularity-dim8}

While singularities are unavoidable in dimension $8$ and higher for specific metrics (as
demonstrated by the Simons cone), a natural and profound question arises:

\begin{question}[Generic Regularity]\label{q:generic-regularity}
Can singularities of minimal hypersurfaces be resolved by small perturbations of the ambient metric
or boundary data?
\end{question}

An affirmative answer to this question, combined with existence results from min-max theory, would
immediately imply the generic existence of smooth minimal hypersurfaces in higher dimensions.

\subsubsection{N. Smale's Pioneering Work}

This question was first addressed by N. Smale \cite{Sma93}, who proved that in any $8$-dimensional
closed manifold with $C^k$-generic metric ($k \geq 4$), there exists a smooth area-minimizing
hypersurface in each nontrivial homology class. Smale's argument crucially relied on the
\emph{Hardt--Simon foliation theorem} \cite{HS85}, which provides a local foliation by smooth
minimal hypersurfaces on either side of a regular area-minimizing cone. The key idea is that
isolated singular points can be ``perturbed away'' by slightly deforming the metric, causing the
singular cone to be replaced by a smooth foliation.

However, Smale's approach requires that the tangent cone at each singular point be one-sided
area-minimizing, a restrictive condition that limits the applicability of the Hardt--Simon local
analysis. Whether regular stable cones that are \emph{not} one-sided area-minimizing admit similar
foliations remains an important open question in the field.

\subsubsection{Recent Breakthroughs: Li--Wang and Beyond}

The generic regularity program has seen spectacular progress in recent years, moving beyond the
limitations of local Hardt--Simon analysis:

\paragraph{Li--Wang: Global Perturbation and Singular Capacity.}
In a groundbreaking work, Yangyang Li and Zhihan Wang \cite{LiWang-dim8} proved that every
$8$-dimensional closed Riemannian manifold with $C^\infty$-generic metric admits a smooth minimal
hypersurface. Their approach represents a paradigm shift in generic regularity theory:

\begin{itemize}
    \item \textbf{Global perturbation framework}: Building on Wang's earlier work \cite{Wang20},
    they developed a global deformation argument that avoids the restrictive assumptions of local
    Hardt--Simon analysis. This allows them to handle singular points whose tangent cones may not be
    one-sided area-minimizing.

    \item \textbf{Singular capacity}: They introduced a novel geometric invariant, the
    \emph{singular capacity}, which counts singular points with suitable weights. This invariant
    decreases under generic perturbations, providing a quantitative measure of progress in resolving
    singularities.

    \item \textbf{Multiplicity considerations}: Their techniques handle minimal hypersurfaces with
    integer multiplicities, addressing a fundamental challenge in combining min-max theory with
    generic regularity.
\end{itemize}

The Li--Wang result settles Question~\ref{q:generic-regularity} affirmatively in dimension $8$ and
demonstrates that singularities in dimension $8$ are indeed ``artifacts'' of special metrics that
disappear under generic perturbations.

\paragraph{Chodosh--Mantoulidis--Schulze--Wang: Second-Order Analysis.}
In a remarkable extension, Chodosh--Mantoulidis--Schulze--Wang \cite{CMSW-dim11} proved generic
regularity up to dimension $11$ for area-minimizers, both in the Plateau problem setting and for
minimizers in homology classes. Their strategy involves a sophisticated second-order blow-up
analysis:

\begin{itemize}
    \item \textbf{Deeper blow-up}: They go one blow-up level deeper than previous works, modeling
    the minimizer locally as a \emph{graph of a Jacobi field} over a cylindrical hypercone
    $\mathcal{C} = \mathcal{C}^\circ \times \mathbb{R}^k$.

    \item \textbf{Almgren frequency at the Jacobi field level}: They introduce a discrete
    Almgren-type frequency monotonicity quantity that measures the decay rate of graph Jacobi fields
    toward the origin. This frequency distinguishes between ``fast decay'' (super-linear) and ``slow
    decay'' (linear) cases.

    \item \textbf{Refined cone-splitting}: In the fast decay case, they obtain improved separation
    estimates between distinct leaves of the foliation. In the slow decay case, new cone-splitting
    arguments at the Jacobi field level provide a second-order refinement of the singular set
    dimension estimates.

    \item \textbf{Non-concentration estimates}: They adapt L. Simon's deep non-concentration
    techniques \cite{Sim93,Sim94} to control the geometry at multiple scales, wrestling with the
    possibility of tangent cone non-uniqueness and the lack of an explicit equation for the singular
    set.
\end{itemize}

Moreover, they proved that in any $\mathbb{R}^{n+1}$, an area-minimizing hypersurface will have
singular set of dimension at most $n - 10 - \varepsilon_n$ after a $C^\infty$-perturbation of the
Plateau boundary, improving the classical $n-7$ bound from Federer's theorem.

\subsection{The Interplay: Min-Max, Integrality, and Regularity}

The regularity theory discussed above is intimately connected to min-max theory. For minimal
hypersurfaces arising specifically from min-max constructions (not just area-minimizers),
Chodosh--Liokumovich--Spolaor \cite{CLS20} proved generic existence of smooth minimal hypersurfaces
in manifolds with positive Ricci curvature. They constructed optimal nested volume-parametrized
sweepouts and combined the non-excessive framework (discussed in Section~\ref{sec:free-bdry}) with
Hardt--Simon type perturbations. This work demonstrates that the non-excessive approach can be
effectively combined with generic regularity techniques, opening new pathways for understanding
singularities in min-max theory.

The connection between integrality, min-max constructions, and regularity forms a central theme of
this thesis. In Chapter~\ref{ch:alternative-regularity}, we develop integrality arguments for varifolds
arising from non-excessive sweepouts when the limiting varifold has multiplicity one. Integrality
is the \emph{crucial bridge} between variational constructions and regularity theory: without it,
Wickramasekera's regularity theorem \cite{Wic14} cannot be applied.

The open problem of establishing integrality for multiplicity $\geq 2$ varifolds in the
non-excessive framework is closely related to the generic regularity questions in higher dimensions.
The techniques of Li--Wang (singular capacity) and Chodosh et al. (second-order blow-up analysis)
suggest potential pathways:

\begin{question}
Can the non-excessive framework be combined with global perturbation techniques (such as the
singular capacity method) to handle higher multiplicity cases and extend generic regularity to
dimension $8$ and beyond in the min-max setting?
\end{question}

Such a combination would unify the variational min-max approach with the recent breakthroughs in
generic regularity theory, potentially creating a comprehensive framework for understanding minimal
hypersurfaces across all dimensions.

\section{The Free Boundary Setting: A New Frontier}\label{sec:free-bdry}

Having discussed the remarkable progress in min-max theory and regularity for \emph{closed}
manifolds (Sections~\ref{sec:minmax}--\ref{sec:regularity}), we now turn to the \emph{free
boundary setting}---the central focus of this thesis. Free boundary minimal hypersurfaces represent
a natural generalization of the closed case, but introduce new geometric and analytic challenges
that require significant technical innovations.

\subsection{From Closed Manifolds to Manifolds with Boundary: Geometric Motivations}

Why study free boundary minimal hypersurfaces? There are several compelling motivations:

\paragraph{Physical motivation:}
In physics, free boundary problems arise naturally when a minimal surface is allowed to ``float
freely'' in a domain, meeting the boundary orthogonally. This orthogonality condition appears in
capillary theory, soap film problems, and extremal surface theory in general relativity.

\paragraph{Geometric motivation:}
From a purely geometric perspective, free boundary minimal hypersurfaces are critical points of the
area functional among all hypersurfaces with boundary lying on $\partial M$. They generalize
geodesics with endpoints on the boundary of a surface, and provide a natural setting for studying
variational problems in manifolds with boundary.

\paragraph{Spectral motivation:}
Free boundary minimal hypersurfaces are intimately connected to the \emph{Steklov eigenvalue
problem}---the spectral problem for the Dirichlet-to-Neumann map on the boundary $\partial M$. This
connection, discussed in Section~\ref{subsec:steklov}, reveals deep relationships between minimal
surface theory, spectral geometry, and boundary value problems.

\paragraph{The challenge:}
Compared to the closed case, the free boundary setting introduces several difficulties:
\begin{itemize}
    \item \textbf{Boundary regularity}: Ensuring that minimal hypersurfaces meet $\partial M$
    smoothly and orthogonally requires refined regularity theory.

    \item \textbf{Compactness}: Varifold compactness is more delicate when boundaries are present,
    as mass can escape to infinity or concentrate at the boundary.

    \item \textbf{Sweepout construction}: Constructing families of hypersurfaces that respect the
    free boundary condition and satisfy the necessary non-triviality conditions is technically
    demanding.
\end{itemize}

Despite these challenges, the free boundary setting has seen significant progress in recent years,
and forms a natural frontier for extending the successes of closed min-max theory.

\subsection{Free Boundary Minimal Hypersurfaces: Definitions and Basic Theory}

\begin{definition}[Free Boundary Minimal Hypersurface]
Let $(M^{n+1}, g)$ be a Riemannian manifold with non-empty boundary $\partial M$. A smooth
hypersurface $\Sigma \subset M$ with boundary $\partial \Sigma \subset \partial M$ is called a
\emph{free boundary minimal hypersurface (FBMH)} if:
\begin{enumerate}
    \item $\Sigma$ is minimal in the interior: $H_\Sigma \equiv 0$ on $\Sigma \setminus \partial \Sigma$,
    where $H_\Sigma$ is the mean curvature of $\Sigma$;

    \item $\Sigma$ meets $\partial M$ orthogonally: $\nu_\Sigma \perp T(\partial M)$ along
    $\partial \Sigma$, where $\nu_\Sigma$ is the unit normal to $\Sigma$.
\end{enumerate}
\end{definition}

The orthogonality condition is the \emph{natural boundary condition} arising from the variational
principle: it ensures that $\Sigma$ is a critical point of the area functional among all
hypersurfaces with boundary on $\partial M$.

\paragraph{Regularity theory for FBMH:}
Li--Zhou \cite{LZ17} developed a comprehensive regularity theory for free boundary minimal
hypersurfaces arising from min-max constructions. They adapted Wickramasekera's regularity theorem
\cite{Wic14} to the free boundary setting, proving that integral stationary varifolds satisfying
the orthogonality condition are smooth up to the boundary in dimensions $3 \leq n+1 \leq 7$. Their
arguments crucially exploit the doubling trick and the almost minimizing property in annuli
intersecting the boundary.

\subsection{The Steklov Eigenvalue Problem}\label{subsec:steklov}

Free boundary minimal hypersurfaces are intimately connected to the \emph{Steklov eigenvalue
problem}, a spectral problem on the boundary $\partial M$. The Steklov operator is defined as the
Dirichlet-to-Neumann map:
\[
    \Lambda: C^\infty(\partial M) \to C^\infty(\partial M), \quad
    \Lambda(f) = \frac{\partial u}{\partial \nu}\Big|_{\partial M},
\]
where $u$ is the harmonic extension of $f$ to $M$ (i.e., $\Delta u = 0$ in $M$ and $u|_{\partial M} = f$).

The Steklov eigenvalues $\{\sigma_k\}_{k=0}^\infty$ are the eigenvalues of $\Lambda$, and the
corresponding eigenfunctions satisfy
\[
    \Delta u = 0 \text{ in } M, \quad
    \frac{\partial u}{\partial \nu} = \sigma u \text{ on } \partial M.
\]

Free boundary minimal hypersurfaces arise as level sets of Steklov eigenfunctions, providing a
variational characterization analogous to the relationship between geodesics and the Laplace
spectrum. This connection has been exploited to study extremal properties of Steklov eigenvalues
and to construct free boundary minimal surfaces in various geometric settings.

\subsection{The Non-Excessive Condition and Integrality}\label{subsec:non-excessive-integrality}

A major recent development in min-max theory is the \emph{non-excessive sweepout framework},
introduced by Chodosh--Liokumovich--Spolaor \cite{CLS20}. This framework provides a conceptual
simplification of the classical Almgren--Pitts approach and has proven particularly effective when
combined with generic regularity techniques.

\subsubsection{What is a Non-Excessive Sweepout?}

A sweepout is called \emph{non-excessive} if it cannot be continuously deformed to reduce its
maximal area while staying in the same homotopy class. This variational criticality condition plays
a role analogous to stationarity, but operates directly at the level of continuous families rather
than through pull-tight constructions.

The non-excessive framework offers both elegance and new challenges. On one hand, it avoids many
technical deformations required in the classical approach. On the other hand, ensuring the
\emph{integrality} of the limiting varifold---a property essential for applying Wickramasekera's
regularity theorem (Section~\ref{subsec:regularity-stable})---requires new arguments.

\subsubsection{The Integrality Challenge}

Integrality is the \emph{crucial bridge} between variational min-max constructions and regularity
theory. Wickramasekera's theorem \cite{Wic14} requires that the stationary varifold be
integral---that is, its multiplicity function takes integer values almost everywhere. Without
integrality, the regularity machinery breaks down, and smoothness cannot be guaranteed.

In the classical Almgren--Pitts framework, integrality is built into the construction through the
use of integral currents and careful pull-tight arguments. However, in the non-excessive framework,
integrality is not automatic and must be established separately. This creates a fundamental
challenge:

\begin{question}[Integrality in the Non-Excessive Framework]
Under what conditions does a stationary varifold arising from a non-excessive sweepout have integer
multiplicities?
\end{question}

This question is addressed in Chapter~\ref{ch:alternative-regularity} of this thesis, where we develop
integrality arguments for the case when the limiting varifold has multiplicity one. Our proof
strategy distinguishes two cases:
\begin{itemize}
    \item \textbf{No mass cancellation}: Standard varifold convergence preserves integrality.

    \item \textbf{With mass cancellation}: We develop new arguments using the non-excessive property
    to control the limiting measure and ensure integer multiplicities.
\end{itemize}

However, a subtle difficulty arises when the limiting varifold has multiplicity two or higher: in
this case, our current techniques do not guarantee integrality, and thus Wickramasekera's regularity
theorem \cite{Wic14} cannot be directly applied. This open problem is deeply connected to the
generic regularity questions discussed in Section~\ref{subsec:generic-regularity-dim8}.

\section{Main Results of This Thesis}

Having laid out the historical and theoretical background, we now state the main results of this
thesis. We present two contributions: (I) the construction of free boundary minimal hypersurfaces
in complete manifolds, and (II) a clarification of the integrality issue in the non-excessive
min-max framework.

\subsection{Result I: Free-Boundary Min-Max in Complete Manifolds}

The primary objective of this work is to establish the existence of free-boundary minimal
hypersurfaces in complete Riemannian manifolds with boundary and sublinear volume growth, combining
the techniques of Chambers--Liokumovich \cite{CL20} for complete manifolds with the free-boundary
theory of Li--Zhou \cite{LZ17}.

\begin{theorem}[Main Theorem: Free Boundary in Complete Manifolds]\label{thm:main}
    Let $3 \leq n + 1 \leq 7$. Let $M^{n + 1}$ be a smooth complete Riemannian manifold with
    non-empty boundary $\partial M$ and sublinear volume growth. Then there exists a hypersurface
    $\Sigma$ with (possibly empty) boundary $\partial\Sigma$ lying on $\partial M$ such that
    \begin{enumerate}
        \item[(i)] $\Sigma$ is smoothly embedded in $M$ up to the boundary $\partial \Sigma$;
        \item[(ii)] $\Sigma$ is minimal in the interior of $\Sigma$; and
        \item[(iii)] $\Sigma$ meets $\partial M$ orthogonally along $\partial \Sigma$.
    \end{enumerate}
    Moreover, $\Sigma$ is smooth up to the boundary.
\end{theorem}

The overall strategy mirrors that of \cite{CL20}, but each step must be revisited to track the
interaction with $\partial M$.

\begin{remark}[Relationship to regularity theory]
Our main theorem (Theorem~\ref{thm:main}) guarantees smoothness in dimensions $3 \leq n+1 \leq 7$
by applying the Li--Zhou free-boundary regularity theory \cite{LZ17} (see
Section~\ref{subsec:lizhou-regularity}). This regularity framework is specifically designed for
free boundary minimal hypersurfaces and operates directly within the min-max construction. The
dimension restriction $3 \leq n+1 \leq 7$ is intrinsic to Li--Zhou's regularity arguments: in
dimension $8$ and higher, singularities may appear even for free boundary minimal hypersurfaces, as
suggested by the existence of the Simons cone in the closed case
(Section~\ref{subsec:generic-regularity-dim8}). The recent breakthroughs of Li--Wang
\cite{LiWang-dim8} and Chodosh et al. \cite{CMSW-dim11} show that generic regularity can be
achieved through metric perturbations in dimensions up to $11$, suggesting potential future
extensions of our complete manifold results to higher dimensions via similar techniques.
\end{remark}

\subsection{Result II: Non-Excessive Min-Max and Integrality}

As a complementary contribution, we provide a new approach to the existence of smooth minimal
hypersurfaces in closed manifolds using the non-excessive sweepout framework, with a focus on
clarifying the \emph{integrality issue}---the crucial bridge between variational constructions and
regularity theory (see Section~\ref{subsec:regularity-stable}).

\subsubsection{Motivation: An Alternative Proof of Pitts' Existence Theorem}

The existence of minimal hypersurfaces in closed Riemannian manifolds was first established by
Pitts~\cite{Pit81} using the Almgren--Pitts min-max theory with regularity provided by
Schoen--Simon and Schoen--Simon--Yau combined with freezing. Subsequently,
Colding--De~Lellis~\cite{CL03} gave an alternative proof for $n+1=3$ (surfaces in 3-manifolds) using
Simon--Smith continuous sweepouts combined with Meeks--Yau regularity, which De~Lellis--Tasnady~\cite{DLT13}
extended to all dimensions $n+1 \leq 7$ using Schoen--Simon--Yau regularity. Both alternative proofs
still use freezing arguments and require working with multiple sweepouts in the construction.

The original motivation for this work was to provide yet another alternative proof of Pitts' theorem
using the non-excessive framework introduced by Chodosh--Liokumovich--Spolaor~\cite{CLS22}. Our
approach differs from previous work in two key aspects: (1) we use Wickramasekera's regularity
theorem~\cite{Wic14} (which requires stability but avoids freezing), with stability provided
directly by the non-excessive condition, and (2) we work with a single, well-structured nested
sweepout rather than the multiple sweepouts required in Colding--De~Lellis and De~Lellis--Tasnady.
We use the same classical tightening procedure as previous approaches to obtain a stationary varifold.

The trade-off is that working with non-excessive sweepouts yields only stationary varifolds, and
integrality must be established separately (whereas it is automatic when working with integral
currents). We succeed in proving integrality for the multiplicity-one case
(Theorem~\ref{thm:non-excessive} below), yielding an alternative proof of Pitts' theorem under this
multiplicity assumption. If one could extend this to multiplicity two, it would yield a complete
alternative proof using Wickramasekera regularity and requiring only a single nested sweepout.

The multiplicity-two case remains open in all dimensions $n \geq 2$, preventing us from fully
realizing the original goal. Nevertheless, our result provides an alternative proof of Pitts'
theorem for multiplicity-one varifolds, offering a different approach to regularity via
Wickramasekera's theorem rather than freezing.

\begin{theorem}[Non-Excessive Min-Max: Multiplicity One Case]\label{thm:non-excessive}
    Let $(M^{n+1}, g)$ be a closed Riemannian manifold with $n \geq 2$. Suppose a non-excessive
    min-max procedure produces a stationary varifold with multiplicity one. Then the varifold is
    integral. Moreover, for $2 \leq n \leq 6$, Wickramasekera's regularity theorem implies the
    varifold is a smooth, closed, embedded minimal hypersurface in $M$.
\end{theorem}

\subsubsection{Key Technical Contributions}

Our main contribution is a \emph{direct proof of integrality} for the stationary varifold arising
from a non-excessive sweepout when it has multiplicity one. As explained in
Section~\ref{subsec:regularity-stable}, integrality is essential for applying Wickramasekera's
regularity theorem \cite{Wic14}: without it, the regularity machinery breaks down and smoothness
cannot be guaranteed.

The proof strategy distinguishes two cases depending on whether mass cancellation occurs in the
limiting process:
\begin{itemize}
    \item \textbf{No mass cancellation}: Standard varifold convergence preserves integrality, and
    we can directly apply Wickramasekera's theorem to conclude smoothness.

    \item \textbf{With mass cancellation}: We develop new arguments using the non-excessive property
    to control the limiting measure. The key insight is that the non-excessive condition prevents
    certain types of mass concentration that would destroy integrality, allowing us to track integer
    multiplicities through the limiting process.
\end{itemize}

\subsubsection{The Multiplicity Barrier and Open Problems}

Our integrality result exhibits a fundamental limitation with respect to multiplicity. For
multiplicity one, our arguments establish integrality in all dimensions $n \geq 2$, and combined
with Wickramasekera's regularity theorem (when $2 \leq n \leq 6$), this yields smooth minimal
hypersurfaces.

However, a crucial difficulty arises when the limiting varifold has multiplicity two or higher. Our
current methods do \emph{not} establish integrality when the limiting varifold has multiplicity
$\geq 2$ in \emph{any} dimension $n \geq 2$. This is the fundamental barrier preventing us from
obtaining a complete alternative proof of Pitts' theorem. The issue is subtle and geometric:

\begin{itemize}
    \item When mass cancellation occurs in the limiting process, two sheets with opposite
    orientations can approach each other, producing a varifold of multiplicity 2 but a current with
    strictly smaller mass.

    \item The multiplicity-one assumption is essential for ruling out such cancellation: if
    $\Theta^n(V,x) = 1$ almost everywhere, cancellation would create points with $\Theta^n \geq 2$,
    yielding a contradiction.

    \item For multiplicity $\geq 2$, this argument breaks down in all dimensions $n \geq 2$, and
    integrality remains an open problem in the non-excessive framework.
\end{itemize}

\textbf{Consequence:} Resolving integrality for multiplicity two would yield a complete alternative
proof of Pitts' existence theorem that is conceptually simpler than both the original
Almgren--Pitts approach and the De~Lellis--Tasnady simplification. However, as it stands, our
result provides only a partial simplification, establishing integrality under the additional
assumption that the limiting varifold has multiplicity one.

This multiplicity barrier is not merely technical---it reflects deep geometric phenomena related to
mass cancellation in the non-excessive framework, and is intimately connected to the generic
regularity questions discussed in Section~\ref{subsec:generic-regularity-dim8}. The recent work of Li--Wang \cite{LiWang-dim8}
introduces a \emph{singular capacity} invariant that counts singular points with suitable weights,
providing a quantitative measure of progress in resolving singularities. Similarly,
Chodosh--Mantoulidis--Schulze--Wang \cite{CMSW-dim11} develop second-order blow-up analysis that
distinguishes between different decay rates of Jacobi fields near singularities.

\begin{question}
Can the non-excessive framework be combined with these global perturbation techniques (such as
singular capacity or second-order blow-up analysis) to handle higher multiplicity cases and extend
generic regularity to dimension $8$ and beyond in the min-max setting?
\end{question}

Such a combination would unify the variational min-max approach with the recent breakthroughs in
generic regularity theory, potentially creating a comprehensive framework for understanding minimal
hypersurfaces that works across all dimensions and multiplicity levels.

\subsection{Organization of the Thesis}

Chapter~\ref{ch:examples} sets the analytic stage for both results: we collect notation, recall
the relative coarea formula, define good sets and good sweepouts, and introduce the non-excessive
framework. The good set framework provides a quantitative way to ensure that width and boundary
size are compatible.

Chapter~\ref{ch:alternative-regularity} presents our non-excessive min-max result for closed manifolds
(Theorem~\ref{thm:non-excessive}). We establish the existence of a stationary varifold via a
simplified tightening argument, then prove its integrality and rectifiability in the multiplicity-one
case. The proof distinguishes two cases depending on whether mass cancellation occurs in the limiting
process. We conclude with a discussion of potential extensions to free-boundary problems and the
challenges posed by higher multiplicities.

Chapters~\ref{ch:structure}--\ref{ch:convergence-of-min-max} contain the proof of the main
free-boundary result (Theorem~\ref{thm:main}) using the continuous sweepout approach.
Chapter~\ref{ch:structure} outlines the logic of the min-max argument and isolates the two major
difficulties: preventing the min-max sequence from escaping to infinity and constructing nested
sweepouts with controlled area. These are resolved in the subsequent technical chapters.
Chapter~\ref{ch:morse-foliation} adapts Morse foliation techniques to the doubled manifold, while
Chapter~\ref{ch:splitting-and-extension} contains relative splitting and extension lemmas needed
to monotonise families. Chapter~\ref{ch:nested-sweepout} uses these tools to build nested sweepouts.
Chapter~\ref{ch:no-escape-to-infinity} proves the non-escape property, and
Chapter~\ref{ch:convergence-of-min-max} carries out the pull-tight and regularity arguments that
culminate in the proof of Theorem~\ref{thm:main}.

Finally, the appendix records the interpolation results that bridge continuous and discrete notions
of almost minimizing, which are used in both the classical and non-excessive approaches.
